\documentclass[pdf,oneside]{oucthesis}
\setmainfont{TeX Gyre Termes}  %设置英文Times New Roman
\setmathfont{TeX Gyre Termes Math} %设置英文数学Times New Roman
% pdf: 便于阅读的电子版，去除掉多余的空白页，加入超链接等
% print：打印版本，添加空白页，便于双面打印，禁用超链接



% ========================================
% 版本控制：电子版 vs 纸质版
% ========================================
\newif\ifprintversion
\printversiontrue   % 纸质版：插入空白页
% \printversionfalse    % 电子版：不插入空白页

% 封皮部分空白页（使用 empty 样式：无页码，无页眉）
\newcommand{\insertblankpageEmpty}{%
    \ifprintversion
        \ifodd\value{page}\else
            \newpage
            \null
            \thispagestyle{empty}
            \newpage
        \fi
    \fi
}


% 前置部分空白页（使用 plain 样式：有页码，无页眉）
\newcommand{\insertblankpagePlain}{%
    \ifprintversion
        \ifodd\value{page}\else
            \newpage
            \null
            \thispagestyle{plain}
            \newpage
        \fi
    \fi
}

% 正文部分空白页（使用 fancy 样式：有页码，有页眉）
\newcommand{\insertblankpageFancy}{%
    \ifprintversion
        \ifodd\value{page}\else
            \newpage
            \null
            \thispagestyle{fancy}
            \newpage
        \fi
    \fi
}


% ========================================
%｜             论文信息填写               ｜
% ========================================
\title{北大西洋亚极地区域多尺度海洋过程能量交换}


% ========================================
%                包引用                 
% ========================================
\usepackage{booktabs}
\usepackage{multirow}
\usepackage{setspace}
\usepackage{amsthm}
\usepackage{algorithm}
\usepackage{algorithmicx}
\usepackage{algpseudocode}
\usepackage{caption}
\captionsetup{font={small}}
\usepackage{csquotes}
\usepackage{adjustbox}
\usepackage{longtable} % 引入 longtable 包，用于跨页表格
\usepackage{array} 
\usepackage{notoccite}  % 忽略目录等部分的引用，解决图表标题中引用参考文献，导致参考文献乱序
\usepackage{seqsplit}  % 解决“长英文导致不分行，超出页面”的问题,使用\seqsplit{}将长英文包括进去即可。
% 内容区域



\raggedbottom %解决latex默认页面底部对齐，导致拉伸出现段间距过大的问题
\begin{document}


% 封面
\include{data/cover}

% 学位论文答辩委员会
\insertblankpageEmpty%空白页判断
\include{data/committee.tex}

% 独创性声明
\insertblankpageEmpty%空白页判断
\include{data/statement.tex}

% 从摘要开始，设置行距为20磅
\setlength{\baselineskip}{20pt}


% 中文摘要
\insertblankpageEmpty%空白页判断
\phantomsection\addengcontent{chapter}{Abstract (in Chinese) }
\include{data/abstract_ch.tex}

\pagenumberingnoreset{Roman}   % 这里将页码设置为大写罗马数值，不然后续插入空白页的时候 页码会变成阿拉伯数字

% 英文摘要
\insertblankpagePlain%空白页判断
\phantomsection\addengcontent{chapter}{Abstract (in English) }
\include{data/abstract_en.tex}

\pagenumberingnoreset{Roman}   % 这里将页码设置为大写罗马数值，不然后续插入空白页的时候 页码会变成阿拉伯数字

% 目录
\insertblankpagePlain%空白页判断
\tableofcontents % 中文目录

\insertblankpagePlain%空白页判断
\tableofengcontents % 英文目录
% 测试文字（用于观察排版位置）

\clearpage
% 图表目录
\insertblankpagePlain%空白页判断
\newcommand{\loflabel}{图} 
\renewcommand{\numberline}[1]{\loflabel~#1\hspace*{1em}}
\renewcommand{\listfigurename}{插~~图~~清~~单} 
% \begin{center}
% \clearpage
% \phantomsection\addcontentsline{toc}{chapter}{插\hspace{6pt}图\hspace{6pt}清\hspace{6pt}单}
% \phantomsection\addengcontent{chapter}{Figure list}
% {\listoffigures}
% \end{center}

% %\clearpage 控制前面目录最后是否要换页，不换页的话，会压紧显示在一页上
\phantomsection
\addcontentsline{toc}{chapter}{插\hspace{6pt}图\hspace{6pt}清\hspace{6pt}单}
\addengcontent{chapter}{Figure list}
\listoffigures
% 上方的注释/显示可以替换

\clearpage
\insertblankpagePlain%空白页判断
\newcommand{\lotlabel}{表}
\renewcommand{\numberline}[1]{\lotlabel~#1\hspace*{1em}}
\renewcommand{\listtablename}{表~~格~~清~~单} 
% \begin{center}
% \clearpage
% \phantomsection\addcontentsline{toc}{chapter}{表\hspace{6pt}格\hspace{6pt}清\hspace{6pt}单}
% \phantomsection\addengcontent{chapter}{Table list}
% {\listoftables}
% \end{center}
% \clearpage
\phantomsection
\addcontentsline{toc}{chapter}{表\hspace{6pt}格\hspace{6pt}清\hspace{6pt}单}
\addengcontent{chapter}{Table list}
\listoftables
% 上方的注释/显示可以替换

\clearpage

% 中文注释表
\insertblankpagePlain%空白页判断
\phantomsection\addengcontent{chapter}{Abbreviation}
\include{data/abbreviation.tex}

% \renewcommand{\numberline}[1]{\lotlabel~#1\hspace*{1em}}
% \renewcommand{\listtablename}{注~~释~~表} 
% \begin{center}
% \clearpage
% \phantomsection\addcontentsline{toc}{chapter}{注\hspace{16pt}释\hspace{16pt}表}
% \phantomsection\addengcontent{chapter}{Abbreviation}
% {\listoftables}
% \end{center}


\pagenumberingnoreset{Roman}   % 这里将页码设置为大写罗马数值，不然后续插入空白页的时候 页码会变成阿拉伯数字
\insertblankpageFancy%空白页判断
% 正文
\mainmatter
\pagenumbering{arabic}

% \mainmatter 会清空页面并重置页码，这里不判断是否是奇数还是偶数页，因为页码是从1开始，转到\mainmatter前判断
\include{contents/section_01}

\insertblankpageFancy%空白页判断
\include{contents/section_02}

\insertblankpageFancy%空白页判断
\include{contents/section_03}

\insertblankpageFancy%空白页判断
\include{contents/section_04}

\insertblankpageFancy%空白页判断
\include{contents/section_05}

\insertblankpageFancy%空白页判断
\include{contents/section_06}

% 参考文献
\insertblankpageFancy%空白页判断
\phantomsection\addengcontent{chapter}{References}
\bibliographystyle{oucauthoryear}
\bibliography{cite}

\clearpage
% 获得的成果
\insertblankpageFancy%空白页判断
\phantomsection\addengcontent{chapter}{Published papers and research results while pursuing the degree}
\include{data/achievement.tex}


% 致谢
\clearpage
\insertblankpageFancy%空白页判断
\phantomsection\addengcontent{chapter}{Acknowledgements}
\include{data/acknowledgement.tex}

% 作者简介
\clearpage
\insertblankpageFancy%空白页判断
\phantomsection\addengcontent{chapter}{Resume}
\include{data/profile.tex}

\end{document}
